\documentclass[12pt,letterpaper]{article}

% ---------- Idioma y codificación ----------
\usepackage[spanish,es-noshorthands]{babel}
\usepackage[utf8]{inputenc}
\usepackage[T1]{fontenc}
\usepackage{csquotes}

% ---------- Estructura y layout ----------
\usepackage{tabularx}
\usepackage{float}
\usepackage{scrhack}
\usepackage[margin=2.5cm]{geometry}

% ---------- Matemáticas ----------
\usepackage{amsmath}

% ---------- Tipografía ----------
\usepackage{newtxtext}
\usepackage{newtxmath}
\usepackage{setspace}
\onehalfspacing

% ---------- Imágenes y tablas ----------
\usepackage{graphicx}
\usepackage{booktabs}
\usepackage{xcolor}
\usepackage{enumitem}
\usepackage[unicode,colorlinks=true,linkcolor=black,citecolor=black,urlcolor=black]{hyperref}
\usepackage{bookmark}

% ---------- Placeholder visual ----------
% Marca el contenido que aún falta completar antes de la entrega.
\newcommand{\relleno}{\textcolor{red}{\itshape [completar]}}

% ---------- Bibliografía ----------
\usepackage[backend=biber, style=apa]{biblatex}
\addbibresource{referencias.bib}

% ---------- Datos que se muestran en el encabezado y pie ----------
\newcommand{\itemUniversity}{Universidad Nacional de San Agustín de Arequipa}
\newcommand{\itemFaculty}{Facultad de Ingeniería de Producción y Servicios}
\newcommand{\itemDepartment}{Departamento Académico de Ingeniería de Sistemas e Informática}
\newcommand{\itemSchool}{Escuela Profesional de Ingeniería de Sistemas}
\newcommand{\itemCourse}{Estadística matemática, probabilidades y métodos empíricos}
\newcommand{\itemStudent}{Grupo N.º 1}

% ---------- Encabezado y pie de página ----------
\usepackage{fancyhdr}
\pagestyle{fancy}
\fancyhf{}
\setlength{\headheight}{40.51407pt}
\renewcommand{\headrulewidth}{1pt}
\renewcommand{\footrulewidth}{1pt}

\fancyhead[L]{\raisebox{-0.2\height}{\includegraphics[width=3cm]{img/logo_episunsa.png}}}
\fancyhead[C]{\fontsize{7}{7}\selectfont \itemUniversity \\ \itemFaculty \\ \itemDepartment \\ \itemSchool \\ \textbf{\itemCourse}}
\fancyhead[R]{\raisebox{-0.2\height}{\includegraphics[width=1.2cm]{img/logo_abet.png}}}

\fancyfoot[L]{\itemStudent}
\fancyfoot[C]{\itemCourse}
\fancyfoot[R]{Página \thepage}

% ============================================================
\begin{document}

% ---------- Portada (archivo externo, sin encabezado ni pie) ----------
% ============================================================
% PORTADA (sin preámbulo: se incluye desde main.tex con % ============================================================
% PORTADA (sin preámbulo: se incluye desde main.tex con % ============================================================
% PORTADA (sin preámbulo: se incluye desde main.tex con \input{portada})
% Datos institucionales fijos. Solo cambia por proyecto:
% título, tipo de trabajo, asignatura, docente y año.
% ============================================================
\begin{titlepage}
\thispagestyle{empty}
\centering
\singlespacing

{\fontsize{18}{21}\selectfont\bfseries Universidad Nacional de San Agustín de Arequipa\par}
\vspace{0.25cm}
{\fontsize{15}{18}\selectfont Facultad de Ingeniería de Producción y Servicios\par}
\vspace{0.22cm}
{\fontsize{15}{18}\selectfont Escuela Profesional de Ingeniería de Sistemas\par}

\vspace{0.9cm}

\includegraphics[width=4.4cm]{img/Escudo_UNSA.png}

\vspace{0.5cm}

{\fontsize{20}{24}\selectfont\bfseries Estimación \par} % Título del trabajo
\vspace{0.5cm}
% {\fontsize{17}{21}\selectfont\bfseries \relleno \par} % Tipo: Monografía / Informe / Glosario / Práctica

\vspace{0.8cm}

{\fontsize{14}{17}\selectfont\textbf{Asignatura:} Estadística matemática, probabilidades y métodos empíricos \par}
\vspace{0.3cm}
{\fontsize{14}{17}\selectfont\textbf{Docente:} Antonia Quispe Mamani \par}

\vspace{0.9cm}

{\fontsize{13}{16}\selectfont\bfseries Integrantes:\par}

\vspace{0.8cm}

\begin{minipage}{0.62\textwidth}
  \raggedright
  \setstretch{1.18}
  {\fontsize{14}{17}\selectfont

        Apaza Chambi Willian Elmer \par
        Capatinta Almanza Camilo Bladimir Alexander \par
        Quispe Hacha Aleyda Luz \par
        Solis Florez Kevin Aldair \par
        Mamani Machaca Paul Roel \par

  }
\end{minipage}

\vfill

{\fontsize{12.5}{15}\selectfont\bfseries Arequipa -- Perú\par}
\vspace{0.2cm}
{\fontsize{12.5}{15}\selectfont\bfseries 2026 \par} % Año o semestre

\end{titlepage})
% Datos institucionales fijos. Solo cambia por proyecto:
% título, tipo de trabajo, asignatura, docente y año.
% ============================================================
\begin{titlepage}
\thispagestyle{empty}
\centering
\singlespacing

{\fontsize{18}{21}\selectfont\bfseries Universidad Nacional de San Agustín de Arequipa\par}
\vspace{0.25cm}
{\fontsize{15}{18}\selectfont Facultad de Ingeniería de Producción y Servicios\par}
\vspace{0.22cm}
{\fontsize{15}{18}\selectfont Escuela Profesional de Ingeniería de Sistemas\par}

\vspace{0.9cm}

\includegraphics[width=4.4cm]{img/Escudo_UNSA.png}

\vspace{0.5cm}

{\fontsize{20}{24}\selectfont\bfseries Estimación \par} % Título del trabajo
\vspace{0.5cm}
% {\fontsize{17}{21}\selectfont\bfseries \relleno \par} % Tipo: Monografía / Informe / Glosario / Práctica

\vspace{0.8cm}

{\fontsize{14}{17}\selectfont\textbf{Asignatura:} Estadística matemática, probabilidades y métodos empíricos \par}
\vspace{0.3cm}
{\fontsize{14}{17}\selectfont\textbf{Docente:} Antonia Quispe Mamani \par}

\vspace{0.9cm}

{\fontsize{13}{16}\selectfont\bfseries Integrantes:\par}

\vspace{0.8cm}

\begin{minipage}{0.62\textwidth}
  \raggedright
  \setstretch{1.18}
  {\fontsize{14}{17}\selectfont

        Apaza Chambi Willian Elmer \par
        Capatinta Almanza Camilo Bladimir Alexander \par
        Quispe Hacha Aleyda Luz \par
        Solis Florez Kevin Aldair \par
        Mamani Machaca Paul Roel \par

  }
\end{minipage}

\vfill

{\fontsize{12.5}{15}\selectfont\bfseries Arequipa -- Perú\par}
\vspace{0.2cm}
{\fontsize{12.5}{15}\selectfont\bfseries 2026 \par} % Año o semestre

\end{titlepage})
% Datos institucionales fijos. Solo cambia por proyecto:
% título, tipo de trabajo, asignatura, docente y año.
% ============================================================
\begin{titlepage}
\thispagestyle{empty}
\centering
\singlespacing

{\fontsize{18}{21}\selectfont\bfseries Universidad Nacional de San Agustín de Arequipa\par}
\vspace{0.25cm}
{\fontsize{15}{18}\selectfont Facultad de Ingeniería de Producción y Servicios\par}
\vspace{0.22cm}
{\fontsize{15}{18}\selectfont Escuela Profesional de Ingeniería de Sistemas\par}

\vspace{0.9cm}

\includegraphics[width=4.4cm]{img/Escudo_UNSA.png}

\vspace{0.5cm}

{\fontsize{20}{24}\selectfont\bfseries Estimación \par} % Título del trabajo
\vspace{0.5cm}
% {\fontsize{17}{21}\selectfont\bfseries \relleno \par} % Tipo: Monografía / Informe / Glosario / Práctica

\vspace{0.8cm}

{\fontsize{14}{17}\selectfont\textbf{Asignatura:} Estadística matemática, probabilidades y métodos empíricos \par}
\vspace{0.3cm}
{\fontsize{14}{17}\selectfont\textbf{Docente:} Antonia Quispe Mamani \par}

\vspace{0.9cm}

{\fontsize{13}{16}\selectfont\bfseries Integrantes:\par}

\vspace{0.8cm}

\begin{minipage}{0.62\textwidth}
  \raggedright
  \setstretch{1.18}
  {\fontsize{14}{17}\selectfont

        Apaza Chambi Willian Elmer \par
        Capatinta Almanza Camilo Bladimir Alexander \par
        Quispe Hacha Aleyda Luz \par
        Solis Florez Kevin Aldair \par
        Mamani Machaca Paul Roel \par

  }
\end{minipage}

\vfill

{\fontsize{12.5}{15}\selectfont\bfseries Arequipa -- Perú\par}
\vspace{0.2cm}
{\fontsize{12.5}{15}\selectfont\bfseries 2026 \par} % Año o semestre

\end{titlepage}

% ---------- Resumen en español ----------
\section{Actividades}

\subsection{Actividad: Estimación}

\subsubsection{Estimación de un parámetro (resumen)}

Un \textbf{parámetro} es una característica numérica de la población (por ejemplo $\mu$, $\sigma^2$ o $P$) que casi siempre es desconocida. La \textbf{estimación} es el proceso de usar los datos de una muestra aleatoria para aproximar ese parámetro. Un \textbf{estimador} es un estadístico (una fórmula) y una \textbf{estimación} es el valor que toma al aplicarlo a una muestra concreta.

Existen dos tipos de estimación:
\begin{itemize}
    \item \textbf{Estimación puntual:} se da un único valor como aproximación del parámetro. Por ejemplo, $\bar{x}$ estima a $\mu$, $s^2$ estima a $\sigma^2$ y $\hat{p}$ estima a $P$. Un buen estimador debe ser insesgado ($E(\hat{\theta})=\theta$), eficiente (de varianza mínima) y consistente (mejora al aumentar $n$).
    \item \textbf{Estimación por intervalo:} se construye un intervalo $[L_1, L_2]$ que contiene al parámetro con una probabilidad $1-\alpha$ llamada \textbf{nivel de confianza}. Tiene la forma general
    \[
        \text{estimador} \pm (\text{valor crítico})\cdot(\text{error estándar}).
    \]
    El nivel de confianza no es la probabilidad de que el parámetro esté en un intervalo ya calculado. Significa que, si se repitiera el muestreo muchas veces, el $(1-\alpha)100\%$ de los intervalos construidos contendría al parámetro.
\end{itemize}

\subsubsection{Problema de aplicación}

El administrador de la red de la facultad quiere evaluar el servidor del aula virtual. El contrato de servicio exige que el servidor cumpla dos condiciones:
\begin{enumerate}
    \item[(i)] El tiempo medio de respuesta no debe superar los $200$ ms.
    \item[(ii)] No más del $10\%$ de las solicitudes deben tardar más de $200$ ms.
\end{enumerate}

Se sabe, por estudios previos, que el tiempo de respuesta se distribuye normalmente, pero se desconocen su media y su varianza. En una hora de trabajo normal se seleccionaron al azar $n=15$ solicitudes y se midió su tiempo de respuesta (en ms):

\begin{center}
\begin{tabular}{|c|c|c|c|c|c|c|c|}
\hline
182 & 195 & 176 & 210 & 188 & 201 & 169 & 193 \\ \hline
185 & 207 & 178 & 191 & 199 & 174 & 186 & \\ \hline
\end{tabular}
\end{center}

Además, en otra muestra aleatoria independiente de $120$ solicitudes, $18$ tardaron más de $200$ ms.

Con un nivel de confianza del $95\%$:
\begin{enumerate}
    \item[a)] Estime por intervalo el tiempo medio de respuesta $\mu$.
    \item[b)] Estime por intervalo la varianza $\sigma^2$ y la desviación estándar $\sigma$.
    \item[c)] Estime por intervalo la proporción $P$ de solicitudes que tardan más de $200$ ms.
    \item[d)] ¿Qué tamaño de muestra se necesita para estimar $\mu$ con un error máximo de $2$ ms? ¿Y para estimar $P$ con un error máximo de $0.05$?
    \item[e)] ¿Cumple el servidor las condiciones del contrato?
\end{enumerate}

\subsubsection{Solución}

\paragraph{Teoría necesaria.}
\begin{enumerate}
    \item \textbf{Estimadores puntuales.} La media muestral $\bar{x}=\frac{1}{n}\sum x_i$ es un estimador insesgado de $\mu$. La varianza muestral $s^2=\frac{1}{n-1}\sum (x_i-\bar{x})^2$ es insesgada para $\sigma^2$. La proporción muestral $\hat{p}=\frac{x}{n}$ es insesgada para $P$.
    \item \textbf{Distribución muestral de la media (Teorema 1).} $\mu_{\bar{x}}=\mu$ y $\sigma_{\bar{x}}^2=\frac{\sigma^2}{n}$. El error estándar de $\bar{x}$ es $\frac{\sigma}{\sqrt{n}}$, que se estima con $\frac{s}{\sqrt{n}}$. Como el muestreo es de una población prácticamente infinita, no se necesita el factor de corrección.
    \item \textbf{Distribución $t$ de Student.} Si la población es normal y $\sigma$ es desconocida, entonces $t=\dfrac{\bar{x}-\mu}{s/\sqrt{n}}$ tiene distribución $t$ con $n-1$ grados de libertad. Por tanto, el intervalo de confianza para $\mu$ es
    \[
        \bar{x} \pm t_{1-\alpha/2,\,n-1}\,\frac{s}{\sqrt{n}}.
    \]
    \item \textbf{Distribución ji-cuadrada.} Si la población es normal, $\chi^2=\dfrac{(n-1)s^2}{\sigma^2}$ tiene distribución $\chi^2$ con $n-1$ grados de libertad. El intervalo de confianza para $\sigma^2$ es
    \[
        \frac{(n-1)s^2}{\chi^2_{1-\alpha/2,\,n-1}} \;\le\; \sigma^2 \;\le\; \frac{(n-1)s^2}{\chi^2_{\alpha/2,\,n-1}}.
    \]
    \item \textbf{Distribución muestral de proporciones.} $\mu_{\hat{p}}=P$ y $\sigma^2_{\hat{p}}=\frac{P(1-P)}{n}$. Si $n\hat{p}>5$ y $n\hat{q}>5$, $\hat{p}$ se aproxima por una normal, y el intervalo es
    \[
        \hat{p} \pm z_{1-\alpha/2}\sqrt{\frac{\hat{p}\,\hat{q}}{n}}.
    \]
    \item \textbf{Tamaño de muestra.} Para un error máximo $E$:
    \[
        n \ge \left(\frac{z_{1-\alpha/2}\,\sigma}{E}\right)^2 \quad \text{(media)}, \qquad
        n \ge \frac{z_{1-\alpha/2}^2\,\hat{p}\,\hat{q}}{E^2} \quad \text{(proporción)}.
    \]
\end{enumerate}

\paragraph{Datos y valores críticos.} Con $1-\alpha=0.95$ se tiene $\alpha/2=0.025$. Del enunciado: $n=15$, $\sum x_i = 2834$.
\[
    \bar{x}=\frac{2834}{15}=188.93\ \text{ms}, \qquad
    s^2=\frac{\sum (x_i-\bar{x})^2}{n-1}=\frac{2054.93}{14}=146.78, \qquad
    s=12.12\ \text{ms}.
\]
Valores de tablas con $n-1=14$ grados de libertad: $t_{0.975,\,14}=2.1448$, $\chi^2_{0.975,\,14}=26.119$, $\chi^2_{0.025,\,14}=5.629$. Para la normal estándar, $z_{0.975}=1.96$.

\paragraph{a) Intervalo para $\mu$.} El error estándar es $\dfrac{s}{\sqrt{n}}=\dfrac{12.12}{\sqrt{15}}=3.13$, de modo que el error de estimación es $2.1448\times 3.13 = 6.71$. Entonces
\[
    188.93 \pm 6.71 \;\Longrightarrow\; 182.22 \le \mu \le 195.64 \ \text{ms}.
\]

\paragraph{b) Intervalo para $\sigma^2$ y $\sigma$.} Con $(n-1)s^2 = 2054.93$:
\[
    \frac{2054.93}{26.119}=78.68 \;\le\; \sigma^2 \;\le\; \frac{2054.93}{5.629}=365.08 .
\]
Tomando raíz cuadrada se obtiene $8.87 \le \sigma \le 19.11$ ms.

\paragraph{c) Intervalo para $P$.} Se tiene $\hat{p}=\dfrac{18}{120}=0.15$ y $\hat{q}=0.85$. Se verifica la condición de normalidad: $n\hat{p}=18>5$ y $n\hat{q}=102>5$. El error estándar es $\sqrt{\dfrac{0.15\times 0.85}{120}}=0.0326$, así que
\[
    0.15 \pm 1.96(0.0326)=0.15 \pm 0.0639 \;\Longrightarrow\; 0.086 \le P \le 0.214 .
\]

\paragraph{d) Tamaño de muestra.}
\begin{itemize}
    \item Para la media, usando $s=12.12$ como aproximación de $\sigma$ y $E=2$:
    \[
        n \ge \left(\frac{1.96\times 12.12}{2}\right)^2 = 140.97 \;\Rightarrow\; n=141 .
    \]
    \item Para la proporción, usando $\hat{p}=0.15$ y $E=0.05$:
    \[
        n \ge \frac{(1.96)^2(0.15)(0.85)}{(0.05)^2}=195.9 \;\Rightarrow\; n=196 .
    \]
    Si no se tuviera información previa, se usaría $\hat{p}=0.5$ (caso más desfavorable) y se necesitarían $n=385$.
\end{itemize}

\paragraph{Resumen de resultados.}
\begin{center}
\begin{tabular}{|l|c|c|}
\hline
\textbf{Parámetro} & \textbf{Estimación puntual} & \textbf{IC al 95\%} \\ \hline
$\mu$ (ms) & $188.93$ & $[182.22;\ 195.64]$ \\ \hline
$\sigma^2$ (ms$^2$) & $146.78$ & $[78.68;\ 365.08]$ \\ \hline
$\sigma$ (ms) & $12.12$ & $[8.87;\ 19.11]$ \\ \hline
$P$ & $0.150$ & $[0.086;\ 0.214]$ \\ \hline
\end{tabular}
\end{center}

\subsubsection{Conclusión}

\begin{itemize}
    \item \textbf{Condición (i), tiempo medio.} Con $95\%$ de confianza, el tiempo medio de respuesta está entre $182.22$ y $195.64$ ms. Todo el intervalo está por debajo de $200$ ms, por lo que hay evidencia suficiente de que el servidor \textbf{cumple} la condición del tiempo medio.
    \item \textbf{Condición (ii), proporción de solicitudes lentas.} El intervalo para $P$ va de $8.6\%$ a $21.4\%$ e incluye el valor $10\%$ exigido. Con esta muestra \textbf{no se puede asegurar} que se cumpla (ni que se incumpla). La estimación puntual ($15\%$) apunta a que podría estar fuera del límite, pero la muestra es insuficiente para concluir.
    \item \textbf{Variabilidad.} El intervalo para $\sigma$ es amplio ($8.87$ a $19.11$ ms), porque con solo $15$ observaciones la varianza se estima con poca precisión. Aunque la media sea buena, una variabilidad alta implica que algunas solicitudes sí superan los $200$ ms, lo que explica el $15\%$ observado.
    \item \textbf{Recomendación.} Para reducir el ancho de los intervalos hay que aumentar la muestra: al menos $141$ observaciones para estimar $\mu$ con error de $2$ ms, y al menos $196$ para estimar $P$ con error de $0.05$. Así se podría decidir con más seguridad si se cumple la condición (ii).
\end{itemize}

Los resultados dependen de que la muestra sea aleatoria y de que el tiempo de respuesta sea aproximadamente normal. Esta última suposición es necesaria para los intervalos de $\mu$ (con $n$ pequeño) y de $\sigma^2$.

\end{document}